\section*{Chapitre \ref{ch:b1:lipids} --- Les lipides}

\begin{solution}{exo:b1:lipids:1}
Un \omterm{def:b1:lipids:fattyacid}{lipide} est une molécule biologique insoluble dans l'eau et soluble
dans les solvants apolaires. \omterm{def:b1:lipids:triglyceride}{Graisses}, \omterm{def:b1:lipids:amphiphilic}{phospholipides}, \omterm{def:b1:lipids:amphiphilic}{stérols} et
caroténoïdes ne partagent aucun squelette commun, seulement ce
comportement vis-à-vis de l'eau, qui leur donne leurs rôles communs
(réserves, films, membranes).
\end{solution}

\begin{solution}{exo:b1:lipids:2}
C18:2 $\Delta^{9,12}$. La dernière double liaison commence au carbone 12
depuis le carboxyle, c'est-à-dire au carbone 6 depuis l'extrémité
méthyle~: oméga-6.
\end{solution}

\begin{solution}{exo:b1:lipids:3}
Un \omterm{def:b1:lipids:triglyceride}{triglycéride} n'a pas de tête \omterm{def:b1:water-small-molecules:water}{polaire}~: entièrement hydrophobe, il est
exclu de l'eau comme phase massive, en gouttelette. Un \omterm{def:b1:lipids:amphiphilic}{phospholipide} est
amphiphile~: sa tête \omterm{def:b1:water-small-molecules:water}{polaire} doit rester dans l'eau alors que ses queues
doivent la quitter, ce que seule une nappe de deux molécules
d'épaisseur peut satisfaire.
\end{solution}

\begin{solution}{exo:b1:lipids:4}
Environ \qty{41}{\celsius}, \qty{-5}{\celsius} et \qty{20}{\celsius}
(la courbe du \omterm{def:b1:lipids:amphiphilic}{cholestérol} n'a pas de transition nette~; la moitié de sa
montée est vers \qty{20}{\celsius}). Une bactérie à \qty{10}{\celsius} a
besoin de la composition insaturée, fluide bien au-dessous de sa
température de culture.
\end{solution}

\begin{solution}{exo:b1:lipids:5}
$15\,000\times 38 = \qty{570}{MJ}$~; $570/8 = 71$ jours. Glycogène~: sec
$570\,000/17 = \qty{33.5}{kg}$, hydraté \qty{134}{kg}.
\end{solution}

\begin{solution}{exo:b1:lipids:6}
C22:0 $>$ C18:0 $>$ C14:0 $>$ C18:1 $>$ C18:3. Parmi les acides saturés,
les chaînes plus longues font plus de contacts de \omterm{def:b1:water-small-molecules:bonds}{van der Waals}~; une
double liaison \emph{cis} supprime plus de tassement que quatre carbones
supplémentaires n'en ajoutent (C18:1 fond au-dessous de C14:0)~; chaque
liaison supplémentaire l'abaisse encore.
\end{solution}

\begin{solution}{exo:b1:lipids:7}
$2\times\qty{140}{\micro m^2}/\qty{0.6}{nm^2} = 2\times 1.4\times
10^{-10}/6\times 10^{-19} = \num{4.7e8}$ molécules~; masse
$\num{4.7e8}\times 750/\num{6e23} = \qty{5.8e-13}{g} = \qty{0.58}{pg}$.
\end{solution}

\begin{solution}{exo:b1:lipids:8}
Volume \qty{1.11}{mL} $= \qty{1.11e-6}{m^3}$~; surface des sphères
$= 6V/d = 6\times 1.11\times 10^{-6}/10^{-6} = \qty{6.7}{m^2}$. Une
goutte unique de ce volume~: $d = \qty{1.28}{cm}$, surface
\qty{5.2}{cm^2}~: l'émulsion offre treize mille fois plus d'interface.
La lipase est hydrosoluble et n'agit qu'à l'interface~; la vitesse y est
proportionnelle.
\end{solution}

\begin{solution}{exo:b1:lipids:9}
Au-dessus de $T_m$, ses cycles rigides gênent le mouvement des chaînes
voisines et réduisent la \omterm{def:b1:lipids:fluidity}{fluidité}~; au-dessous de $T_m$, son
encombrement empêche les chaînes de se tasser en gel ordonné et empêche
le gel. Aucun tassement coopératif n'étant possible, il n'y a pas de
température à laquelle toute la bicouche change d'état d'un coup~: la
transition est étalée.
\end{solution}

\begin{solution}{exo:b1:lipids:10}
Chaînes toutes saturées~: à \qty{25}{\celsius} les membranes sont déjà
proches de leur transition et raides~; à \qty{5}{\celsius} elles
gélifient. Les membranes gélifiées laissent les protéines cesser de
fonctionner et, au réchauffement ou sous une contrainte mécanique, se
fissurent et fuient~: les \omterm{def:b1:cell-unit-of-life:cell}{cellules} perdent leur contenu et le \omterm{def:b1:body-plans-tissues:tissue}{tissu}
meurt (dommage au froid). Le type sauvage désature ses \omterm{def:b1:lipids:fattyacid}{lipides} à
l'automne et reste fluide à \qty{5}{\celsius}.
\end{solution}

\begin{solution}{exo:b1:lipids:11}
\qty{1.07}{kg} d'eau par kilogramme de \omterm{def:b1:lipids:triglyceride}{graisse}. Oxygène~: \qty{2000}{L}~;
à \qty{5}{\%} d'extraction sur \qty{21}{\%} d'oxygène, l'air nécessaire
vaut $2000/(0.21\times 0.05) = \qty{190}{m^3}$~; l'eau perdue à
\qty{30}{g/m^3}~: \qty{5.7}{kg}. Le chameau perd cinq fois plus d'eau en
respirant que la \omterm{def:b1:lipids:triglyceride}{graisse} n'en fournit~: la bosse est une réserve
d'énergie, non une réserve d'eau.
\end{solution}

\begin{solution}{exo:b1:lipids:12}
Une bicouche étant donnée, l'\omterm{prop:b1:water-small-molecules:properties}{effet hydrophobe} la fait se recoller, se
refermer en vésicules et accepter de nouveaux \omterm{def:b1:lipids:fattyacid}{lipides} sans apport
d'énergie~: l'assemblage est spontané. Mais la \omterm{def:b1:cell-unit-of-life:cell}{cellule} fabrique ses
\omterm{def:b1:lipids:fattyacid}{lipides} avec des enzymes logées dans une membrane préexistante (le RE),
qui les y insèrent sur place~; une membrane neuve croît par expansion
d'une ancienne et est répartie à la division, de sorte que toute
membrane de toute \omterm{def:b1:cell-unit-of-life:cell}{cellule} descend des membranes de la \omterm{def:b1:cell-unit-of-life:cell}{cellule}
précédente. L'auto-assemblage explique la physique de la nappe, non son
origine dans la \omterm{def:b1:cell-unit-of-life:cell}{cellule}.
\end{solution}

\begin{solution}{pb:b1:lipids:1}
\textbf{1.} C16:0 (0), C18:1 $\Delta^9$ (1), C22:6 $\Delta^{4,7,10,13,16,19}$
(6).
\textbf{2.} Truite chaude~: $0.40\times 0 + 0.45\times 1 + 0.15\times 6 = 1.35$~;
truite froide~: $0 + 0.40 + 2.10 = 2.50$ doubles liaisons par chaîne.
\textbf{3.} Chaque liaison \emph{cis} coude la chaîne et empêche le
tassement étroit avec les voisines~; moins de contacts de \omterm{def:b1:water-small-molecules:bonds}{van der Waals},
moins de chaleur nécessaire pour désordonner les chaînes, $T_m$ plus
bas.
\textbf{4.} $1.15$ liaison de plus par chaîne~: environ
\qty{45}{\celsius} de moins.
\textbf{5.} Dans un gel, les \omterm{def:b1:lipids:fattyacid}{lipides} ne peuvent plus bouger~: les
transporteurs ne peuvent plus changer de conformation, les \omterm{def:b1:membranes-transport:transporters}{pompes}
s'arrêtent, les canaux se bloquent~; les gradients ioniques
s'épuisent, la conduction nerveuse échoue, et le poisson est paralysé
puis meurt d'un froid qu'un poisson acclimaté tolère.
\textbf{6.} Les désaturases, qui introduisent des doubles liaisons dans
les chaînes d'acyles gras~; ce sont des \omterm{def:b1:membranes-transport:proteins}{protéines intrinsèques} du
\omterm{def:b1:eukaryotic-cell:endomembrane}{réticulum endoplasmique}, où les \omterm{def:b1:lipids:fattyacid}{lipides} sont fabriqués.
\textbf{7.} Dans les deux cas il tamponne la \omterm{def:b1:lipids:fluidity}{fluidité}~: il raidit la
membrane froide très insaturée au-dessus de son bas $T_m$, et il empêche
la membrane chaude de gélifier lors d'une journée fraîche.
\textbf{8.} $30\,000\times 38 = \qty{1140}{MJ}$.
\textbf{9.} $1140/60 = 19$ jours.
\textbf{10.} Sec~: $\num{1.14e6}/17 = \qty{67}{kg}$~; hydraté~:
\qty{268}{kg}.
\textbf{11.} $268 - 30 = \qty{238}{kg}$, près de la moitié de la masse
du chameau.
\textbf{12.} Une couche de \omterm{def:b1:lipids:triglyceride}{graisse} sous toute la peau isolerait le corps
et l'empêcherait d'évacuer sa chaleur dans le désert~; concentrer la
\omterm{def:b1:lipids:triglyceride}{graisse} en une seule bosse laisse le reste de la peau libre de perdre de
la chaleur.
\textbf{13.} $1.07\times 30 = \qty{32}{kg}$ d'eau.
\textbf{14.} $2.0\times 30\,000 = \qty{60000}{L} = \qty{60}{m^3}$
d'oxygène.
\textbf{15.} $60/(0.21\times 0.05) = \qty{5700}{m^3}$ d'air.
\textbf{16.} $5700\times 30 = \qty{171}{kg}$ d'eau.
\textbf{17.} Cinq fois plus perdu que gagné~: oxyder la bosse coûte de
l'eau. La bosse est une réserve d'énergie qui permet au chameau de se
passer de nourriture~; son économie d'eau vient de ses reins, de sa
tolérance à la déshydratation et de sa capacité à laisser monter sa
température.
\textbf{18.} $500\times 3.5\times 6 = \qty{10.5}{MJ}$ stockés en chaleur et libérés la nuit~; les évacuer par évaporation aurait coûté $10\,500/2.4 = \qty{4.4}{kg}$ d'eau par jour.
\textbf{19.} Une membrane ne peut s'étirer que de quelques pour cent~;
doubler le volume d'un globule rouge exige donc un large excédent de
membrane replié dans la forme biconcave de repos, et un \omterm{def:b1:eukaryotic-cell:cytoskeleton}{cytosquelette}
qui la laisse se déplier sans se déchirer~: les globules rouges du
chameau sont ovales, avec plus de membrane par unité de volume que ceux
d'un humain.
\textbf{20.} $2\times 80\times 10^{-12}/0.65\times 10^{-18} = \num{2.5e8}$.
\textbf{21.} $\num{2.5e8}\times\num{8e12}\times 40 = \num{7.9e22}$.
\textbf{22.} $\num{7.9e22}\times 750/\num{6e23} = \qty{99}{g}$.
\textbf{23.} La diffusion latérale garde à chaque instant la tête dans
l'eau et les queues dans l'\omterm{def:b1:lipids:triglyceride}{huile}, ce qui ne coûte rien~; la bascule
exige que la tête \omterm{def:b1:water-small-molecules:water}{polaire} traverse le cœur hydrophobe, une grande
barrière d'énergie que l'agitation thermique ne franchit qu'une fois par
semaine.
\textbf{24.} Une différence entre feuillets ne s'efface qu'à la vitesse
à laquelle les \omterm{def:b1:lipids:fattyacid}{lipides} basculent~; à une fois par semaine et par
molécule, une asymétrie établie par des enzymes qui font traverser des
\omterm{def:b1:lipids:fattyacid}{lipides} précis (les flippases) persiste indéfiniment contre la bascule
spontanée.
\textbf{25.} Environ \qty{240}{kg} économisés --- près de la moitié de
sa masse corporelle --- en portant \qty{30}{kg} de \omterm{def:b1:lipids:triglyceride}{graisse} au lieu de
\qty{270}{kg} de glycogène hydraté~; et la bosse n'est pas une réserve
d'eau~: l'oxyder fait perdre dans le souffle cinq fois plus d'eau qu'elle
n'en produit.
\end{solution}
